% id: 0505
% book: RPM 확률과 통계 · 05 확률분포 (2)
% section: 실력 Up
% figure: crop:fig-0505.png
% answer: ①
% answer_source: 답지
% variant_level: 0 (원본 전사)
% 이 파일은 build.mjs 가 items.json 에서 만든다. 고칠 때는 items.json 을 고친다.
\smash{\raisebox{1.5mm}{\dmkichul{RPM 확률과 통계 0505번 · 84쪽 · 교육청 기출}}}
\noindent\begin{minipage}[t]{0.58\linewidth}\vspace{0pt}\raggedright 정규분포를 따르는 두 확률변수 $X$, $Y$의 확률밀도함수는 각각 $f(x)$, $g(x)$이다. $\mathrm{V}(X)=\mathrm{V}(Y)$이고, 양수 $a$에 대하여 \par\smallskip\end{minipage}\hfill\begin{minipage}[t]{0.4\linewidth}\vspace{0pt}\centering \setlength{\dmfigw}{106.7pt}\ifdim\dmfigw>\linewidth\setlength{\dmfigw}{\linewidth}\fi\includegraphics[width=\dmfigw]{figures/fig-0505.png}\end{minipage}\par
\noindent \dispeq{$f(a)=f(3a)=g(2a)$,}\par\dispeq{$\mathrm{P}(Y\le 2a)=0.6915$}\par\smallskip\noindent 일 때, $\mathrm{P}(0\le X\le 3a)$의 값을 오른쪽 표준정규분포표를 이용하여 구한 것은?
\begin{choices32}{$0.5328$}{$0.6247$}{$0.6687$}{$0.7745$}{$0.8185$}\end{choices32}
